% GENERATED FILE. Edit canonical lessons, then run npm run generate:books.
% canonical-source-hash: fnv1a64:2edb77c5dd2e5694
% canonical-lessons: MW-C284-basto, MW-C284-pati, MW-C284-syahi, MW-C284-davat, MW-C284-bhanno, MW-R284-review-279-284

\chapter{My School Bag}
\label{ch:mw-my-school-bag}
\hlchaptermodality{284}

\begin{chapteropening}
\textbf{By the end of this chapter:} \emph{I can say what is in my school bag and what I study.}

Complete the last lesson of chapter 284: I can say what is in my school bag and what I study.
\end{chapteropening}

\section[\texorpdfstring{\mw{बस्तो} (basto)}{बस्तो (basto)}]{\mw{बस्तो} (basto) — a school bag}

\label{lesson:MW-C284-basto}



\begin{quote}
Before the new one: say the Marwadi for embroidery, then the Marwadi for to decorate.
\end{quote}

\subsection*{You'll want to know: \mw{बस्तो}}
\textbf{\mw{बस्तो}} — \emph{basto} — \textquotedblleft{}a school bag\textquotedblright{}.

\begin{quote}
\mw{म्हारो} \mw{बस्तो} \mw{भारी} \mw{है।}
\end{quote}

\begin{quote}
\emph{mhāro basto bhārī hai.}
\end{quote}

\begin{quote}
My school bag is heavy.
\end{quote}

A masculine noun.

\begin{cousinweb}
\textbf{\mw{बस्तो}} (\emph{basto}) is usually traced to Persian \emph{basta}, \textquotedblleft{}a bundle\textquotedblright{}; Marwadi gives it the -\mw{ो} ending of a masculine noun.
\end{cousinweb}

\begin{grammarlens}[title={What you've built}]
One more word for this chapter.
\end{grammarlens}

\subsection*{Your turn}
\begin{itemize}
\raggedright
  \item \emph{Say it:} \emph{basto}
  \item \emph{Say it:} \emph{basto}, once more
  \item \emph{Say it:} \emph{mhāro basto bhārī hai.} — My school bag is heavy.
  \item \emph{Recall:} say the Marwadi for embroidery, then the Marwadi for to decorate, then say \emph{basto} again
\end{itemize}

\begin{tcolorbox}[breakable,colback=teal!4,colframe=teal!35!black,title={Before you move on}]
What does \mw{बस्तो} mean? (\textquotedblleft{}A school bag\textquotedblright{}.) Say it once more.
\end{tcolorbox}
\section[\texorpdfstring{\mw{पाटी} (pāṭī)}{पाटी (pāṭī)}]{\mw{पाटी} (pāṭī) — a slate}

\label{lesson:MW-C284-pati}



\begin{quote}
Before the new one: say the Marwadi for to decorate, then the Marwadi for a school bag.
\end{quote}

\subsection*{You'll want to know: \mw{पाटी}}
\textbf{\mw{पाटी}} — \emph{pāṭī} — \textquotedblleft{}a slate\textquotedblright{}.

\begin{quote}
\mw{म्हारी} \mw{पाटी} \mw{काळी} \mw{है।}
\end{quote}

\begin{quote}
\emph{mhārī pāṭī kāḷī hai.}
\end{quote}

\begin{quote}
My slate is black.
\end{quote}

A feminine noun.

\begin{cousinweb}
\textbf{\mw{पाटी}} (\emph{pāṭī}) is inherited from Sanskrit \emph{paṭṭikā}, \textquotedblleft{}a slab, a tablet\textquotedblright{}; CDIAL lists Old Marwari \emph{pāṭī}, \textquotedblleft{}plank\textquotedblright{}, and Kumaoni \emph{pāṭī}, \textquotedblleft{}board, writing board\textquotedblright{} (CDIAL 7699).
\end{cousinweb}

\begin{grammarlens}[title={What you've built}]
One more word for this chapter.
\end{grammarlens}

\subsection*{Your turn}
\begin{itemize}
\raggedright
  \item \emph{Say it:} \emph{pāṭī}
  \item \emph{Say it:} \emph{pāṭī}, once more
  \item \emph{Say it:} \emph{mhārī pāṭī kāḷī hai.} — My slate is black.
  \item \emph{Recall:} say the Marwadi for to decorate, then the Marwadi for a school bag, then say \emph{pāṭī} again
\end{itemize}

\begin{tcolorbox}[breakable,colback=teal!4,colframe=teal!35!black,title={Before you move on}]
What does \mw{पाटी} mean? (\textquotedblleft{}A slate\textquotedblright{}.) Say it once more.
\end{tcolorbox}
\section[\texorpdfstring{\mw{स्याही} (syāhī)}{स्याही (syāhī)}]{\mw{स्याही} (syāhī) — ink}

\label{lesson:MW-C284-syahi}



\begin{quote}
Before the new one: say the Marwadi for a school bag, then the Marwadi for a slate.
\end{quote}

\subsection*{You'll want to know: \mw{स्याही}}
\textbf{\mw{स्याही}} — \emph{syāhī} — \textquotedblleft{}ink\textquotedblright{}.

\begin{quote}
\mw{स्याही} \mw{नीली} \mw{है।}
\end{quote}

\begin{quote}
\emph{syāhī nīlī hai.}
\end{quote}

\begin{quote}
The ink is blue.
\end{quote}

A feminine noun.

\begin{cousinweb}
\textbf{\mw{स्याही}} (\emph{syāhī}) is usually traced to Persian \emph{siyāhī}, \textquotedblleft{}blackness\textquotedblright{}, and so \textquotedblleft{}ink\textquotedblright{}.
\end{cousinweb}

\begin{grammarlens}[title={What you've built}]
One more word for this chapter.
\end{grammarlens}

\subsection*{Your turn}
\begin{itemize}
\raggedright
  \item \emph{Say it:} \emph{syāhī}
  \item \emph{Say it:} \emph{syāhī}, once more
  \item \emph{Say it:} \emph{syāhī nīlī hai.} — The ink is blue.
  \item \emph{Recall:} say the Marwadi for a school bag, then the Marwadi for a slate, then say \emph{syāhī} again
\end{itemize}

\begin{tcolorbox}[breakable,colback=teal!4,colframe=teal!35!black,title={Before you move on}]
What does \mw{स्याही} mean? (\textquotedblleft{}Ink\textquotedblright{}.) Say it once more.
\end{tcolorbox}
\section[\texorpdfstring{\mw{दवात} (davāt)}{दवात (davāt)}]{\mw{दवात} (davāt) — an inkpot}

\label{lesson:MW-C284-davat}



\begin{quote}
Before the new one: say the Marwadi for a slate, then the Marwadi for ink.
\end{quote}

\subsection*{You'll want to know: \mw{दवात}}
\textbf{\mw{दवात}} — \emph{davāt} — \textquotedblleft{}an inkpot\textquotedblright{}.

\begin{quote}
\mw{दवात} \mw{मेज} \mw{माथै} \mw{है।}
\end{quote}

\begin{quote}
\emph{davāt mej māthai hai.}
\end{quote}

\begin{quote}
The inkpot is on the table.
\end{quote}

A feminine noun.

\begin{cousinweb}
\textbf{\mw{दवात}} (\emph{davāt}) is usually traced, through Persian \emph{davāt}, to Arabic \emph{dawāt}, \textquotedblleft{}an inkstand\textquotedblright{}.
\end{cousinweb}

\begin{grammarlens}[title={What you've built}]
One more word for this chapter.
\end{grammarlens}

\subsection*{Your turn}
\begin{itemize}
\raggedright
  \item \emph{Say it:} \emph{davāt}
  \item \emph{Say it:} \emph{davāt}, once more
  \item \emph{Say it:} \emph{davāt mej māthai hai.} — The inkpot is on the table.
  \item \emph{Recall:} say the Marwadi for a slate, then the Marwadi for ink, then say \emph{davāt} again
\end{itemize}

\begin{tcolorbox}[breakable,colback=teal!4,colframe=teal!35!black,title={Before you move on}]
What does \mw{दवात} mean? (\textquotedblleft{}An inkpot\textquotedblright{}.) Say it once more.
\end{tcolorbox}
\section[\texorpdfstring{\mw{भणणो} (bhaṇṇo)}{भणणो (bhaṇṇo)}]{\mw{भणणो} (bhaṇṇo) — to study}

\label{lesson:MW-C284-bhanno}



\begin{quote}
Before the new one: say the Marwadi for ink, then the Marwadi for an inkpot.
\end{quote}

\subsection*{You'll want to know: \mw{भणणो}}
\textbf{\mw{भणणो}} — \emph{bhaṇṇo} — \textquotedblleft{}to study\textquotedblright{}.

\begin{quote}
\mw{हूं} \mw{स्कूल} \mw{मांय} \mw{भणूं} \mw{हूं।}
\end{quote}

\begin{quote}
\emph{hū̃ skūl mā̃y bhaṇū̃ hū̃.}
\end{quote}

\begin{quote}
I study at school.
\end{quote}

Past (I …): \mw{म्हैं} \mw{किताब} \mw{भणी}, \emph{mhaĩ kitāb bhaṇī}.

\begin{cousinweb}
\textbf{\mw{भणणो}} (\emph{bhaṇṇo}) is inherited from Sanskrit \emph{bhaṇati}, \textquotedblleft{}calls aloud, speaks\textquotedblright{}; CDIAL lists Old Marwari \emph{bhaṇaï}, \textquotedblleft{}recites\textquotedblright{}, and Gujarati \emph{bhaṇvũ}, \textquotedblleft{}to say, repeat, study\textquotedblright{} (CDIAL 9383). It is the Marwadi word for studying.
\end{cousinweb}

\begin{grammarlens}[title={What you've built}]
One more word for this chapter.
\end{grammarlens}

\subsection*{Your turn}
\begin{itemize}
\raggedright
  \item \emph{Say it:} \emph{bhaṇṇo}
  \item \emph{Say it:} \emph{bhaṇṇo}, once more
  \item \emph{Say it:} \emph{hū̃ skūl mā̃y bhaṇū̃ hū̃.} — I study at school.
  \item \emph{Recall:} say the Marwadi for ink, then the Marwadi for an inkpot, then say \emph{bhaṇṇo} again
\end{itemize}

\begin{tcolorbox}[breakable,colback=teal!4,colframe=teal!35!black,title={Before you move on}]
What does \mw{भणणो} mean? (\textquotedblleft{}To study\textquotedblright{}.) Say it once more.
\end{tcolorbox}
\section[(dialogue)]{Looking back: from Games and Contests to My School Bag}

\label{lesson:MW-R284-review-279-284}



\begin{quote}
Say the Marwadi for a friend and a board game.
\end{quote}

\subsection*{Your turn}
Say these aloud:

\begin{itemize}
\raggedright
  \item a friend, a board game, a die for games, a contest, wrestling
  \item a team, a ball, to catch a ball, a target, to take aim
  \item a potter's wheel, to spin a wheel, to pat into shape, to carve, a kiln
  \item yarn, a spinning wheel, to spin thread, a garland, to string beads
  \item a pearl, to pierce, to inlay, embroidery, to decorate
  \item a school bag, a slate, ink, an inkpot, to study
\end{itemize}

\begin{tcolorbox}[breakable,colback=teal!4,colframe=teal!35!black,title={Before you move on}]
A friend? (\textbf{\mw{गोठियो}}, \emph{goṭhiyo}.) A board game? (\textbf{\mw{चौपड़}}, \emph{chaupaṛ}.) A die for games? (\textbf{\mw{पासो}}, \emph{pāso}.) A contest? (\textbf{\mw{होड}}, \emph{hoḍ}.) Wrestling? (\textbf{\mw{कुस्ती}}, \emph{kustī}.) A team? (\textbf{\mw{टीम}}, \emph{ṭīm}.) A ball? (\textbf{\mw{गेंद}}, \emph{gẽd}.) To catch a ball? (\textbf{\mw{झेलणो}}, \emph{jhelṇo}.) A target? (\textbf{\mw{निसाणो}}, \emph{nisāṇo}.) To take aim? (\textbf{\mw{सांधणो}}, \emph{sā̃dhṇo}.) A potter's wheel? (\textbf{\mw{चाक}}, \emph{chāk}.) To spin a wheel? (\textbf{\mw{घुमावणो}}, \emph{ghumāvṇo}.) To pat into shape? (\textbf{\mw{थापणो}}, \emph{thāpṇo}.) To carve? (\textbf{\mw{कोरणो}}, \emph{korṇo}.) A kiln? (\textbf{\mw{भट्टी}}, \emph{bhaṭṭī}.) Yarn? (\textbf{\mw{सूत}}, \emph{sūt}.) A spinning wheel? (\textbf{\mw{चरखो}}, \emph{charkho}.) To spin thread? (\textbf{\mw{कातणो}}, \emph{kātṇo}.) A garland? (\textbf{\mw{माळा}}, \emph{māḷā}.) To string beads? (\textbf{\mw{गूंथणो}}, \emph{gū̃thṇo}.) A pearl? (\textbf{\mw{मोती}}, \emph{motī}.) To pierce? (\textbf{\mw{बींधणो}}, \emph{bī̃dhṇo}.) To inlay? (\textbf{\mw{जड़णो}}, \emph{jaṛṇo}.) Embroidery? (\textbf{\mw{कसीदो}}, \emph{kasīdo}.) To decorate? (\textbf{\mw{सजावणो}}, \emph{sajāvṇo}.) A school bag? (\textbf{\mw{बस्तो}}, \emph{basto}.) A slate? (\textbf{\mw{पाटी}}, \emph{pāṭī}.) Ink? (\textbf{\mw{स्याही}}, \emph{syāhī}.) An inkpot? (\textbf{\mw{दवात}}, \emph{davāt}.) To study? (\textbf{\mw{भणणो}}, \emph{bhaṇṇo}.)
\end{tcolorbox}
